
\documentclass[11pt]{article}

\usepackage{amsfonts}
\usepackage{amsmath}
\usepackage{amscd}
\usepackage{amssymb}
\usepackage{indentfirst}
\usepackage{natbib}
\usepackage[colorlinks=true,allcolors=blue]{hyperref}
\usepackage{url}
\usepackage{doi}
\usepackage{tikz-cd}

\DeclareMathOperator{\idfun}{id}
\DeclareMathOperator{\isContr}{isContr}
\DeclareMathOperator{\isProp}{isProp}
\DeclareMathOperator{\isSet}{isSet}
\DeclareMathOperator{\dom}{dom}
\DeclareMathOperator{\cod}{cod}
\DeclareMathOperator{\qinv}{qinv}
\DeclareMathOperator{\binv}{biinv}
\DeclareMathOperator{\isEquiv}{isEquiv}
\DeclareMathOperator{\refl}{refl}

\newcommand{\set}[1]{\{\,#1\,\}}

\newcommand{\nats}{\mathbb{N}}
\newcommand{\real}{\mathbb{R}}

\newcommand{\fatdot}{\,\cdot\,}

\newcommand{\REVISED}{\begin{center} \LARGE REVISED DOWN TO HERE \end{center}}
\newcommand{\MOVED}[1][equation]{\begin{center} [#1 moved] \end{center}}

\begin{document}

\title{On Generalizations of Equality, Logic, and Set Theory
    in Mathematics of the Past Century}
\author{Charles J. Geyer}
\maketitle

\section{Introduction}

This document started with two questions raised by Roy Cook
in the foundations interest
group (FIG) of the Minnesota Center for Philosophy of Science, one March 19,
2026 and one much earlier (back in fall 2025).
\begin{itemize}
\item What is really a substructure?
\item What is finite, and can you define the natural
    numbers when you don't even know what finite is?
\end{itemize}

The point of this document is to explain how contemporary mathematics
makes these
questions uninteresting.  They do not even arise.  Those brainwashed in
logic and traditional
set theory cannot understand what this is about without learning
newer ideas from abstract algebra, category theory, constructive analysis,
and type theory, especially homotopy type theory.

So this document has evolved into a history of mathematics since 1900
with special attention to these subjects.

\section{Substructure}

\subsection{Up to Isomorphism}

Long ago in FIG we read \citet{mazur}, which explains that ``equal''
in contemporary mathematics has been demoted from logical equality
to ``isomorphic.''

How many trivial groups (having exactly one element) are there?
The up-to-isomorphism view says just one.
If you insist that there are many, then all are isomorphic.
Contemporary mathematics says that arguing about whether the only element of a
trivial group is ``really'' the empty set or a banana or Julius Caesar
is just silly.

And the same for everything else in mathematics, equality up to isomorphism
is good enough.
I know about Frege's Caesar problem.  But I don't care, and neither does
any contemporary mathematician.
Of course, even structuralist philosophers agree with this position, more
or less (of course there are many different structuralist views).

\subsection{The Elementary Theory of the Category of Sets}

Set theory was reconceptualized in the 1960's by
Lawvere.  His elementary theory of the category of sets (ETCS)
\citep{lawvere-etcs}
has a different way of making the issue of substructure disappear.

Long ago in FIG we read \citet{leinster}, which is an elementary
explication of ETCS.  ETCS is motivated by category theory and topos theory,
but is actually just a theory in first-order logic like ZFC
(Zermelo–Fraenkel set theory with the axiom of choice).  ZFC is a one-sorted
theory with variables ranging over sets and one non-logical binary relation
symbol $\in$ read as membership.
ETCS is a two-sorted
theory with some variables ranging over sets and some variables ranging over
functions and one non-logical binary relation
symbol $\circ$ read as composition of functions.
We can (this is a matter of taste) also include non-logical constants
$\idfun_A$ for every set $A$, which are identity functions.
If we do not include these constants,
instead we include axioms that say identity functions exist and are unique.

In short, in ZFC sets and set membership are the primary concepts,
whereas in ETCS sets and
functions and function composition are the primary concepts, and the
set membership and subset relations must be defined in terms of them.

To define membership \citep[Section~1]{leinster}
we first define \emph{terminal} sets.
These are sets $T$ such there is at exactly one function $A \to T$ for any
set $A$.  An axiom of ETCS asserts that terminal sets exist.
Pick one and denote it $1$ (it turns out that, as in category theory,
all terminal sets are isomorphic,\footnote{Proof: Suppose $T$ and $T'$ are
terminal.  Then the functions $T \to T' \to T$ and $T' \to T \to T'$ are
unique, hence must be the identity functions $\idfun_T$ and $\idfun_{T'}$.
But this says the unique functions $T \to T'$ and $T' \to T$ are
inverse functions of each other, hence isomorphisms.
\label{foot:terminal}
See Section~\ref{sec:category} below for the definition of isomorphism.}
so it does not matter which one we pick).

Then ETCS defines an element $x$ of a set $A$ to be a function $x : 1 \to A$.
So worrying about whether (or exactly how) $x$ is in $A$ is a non-starter.
Functions (including $x$) aren't in sets, rather they go between sets
(from one to another).

You may be wondering what function evaluation now means.  What does $f(x)$
denote if $x$ is not an element of $A$ in the sense of ZFC?  It denotes
the function composition $f \circ x$, that is, the composition of functions
$$
\begin{CD}
   1 @>x>> A @>f>> B
\end{CD}
$$

This all makes sense because the set $1$ has exactly one element
$* : 1 \to 1$ (unique because 1 is terminal).  So for any $x$ as above,
$f(x)$ means the same as $f(x(*))$.  And this is an element of the codomain
of $f$.

To define the subset relation \citep[Discussion of Axiom~8]{leinster}
we first define \emph{injective} functions.
In ETCS we say a function is \emph{injective} if
it has the left cancellation property: if $x$ and $y$ are elements of
the domain of $f$, then $f \circ x = f \circ y$ implies $x = y$.
Of course, this is just the usual definition in different notation.

In ETCS a \emph{subset} is just an injective function into a set.
Hence it makes no philosophical sense to ask in what way this function
has some sort of part-whole relationship to its codomain.

\subsection{Abstract Algebra}

The up-to-isomorphism view goes back to the rise of abstract algebra in
the 1920's in the work of Noether and coworkers
\citep[Chapter~6]{kleiner} and in the textbooks of
\citet{van-der-waerden} and \citet{birkhoff-maclane}.

Consider abstract group theory, so the only sets of interest are groups
and the only functions of interest are group homomorphisms.
The first isomorphism theorem (of group theory,
\citealp[Theorem~8.1]{aluffi}, almost every algebraic theory
has a similar theorem) says, if $f : G \to H$ is a group homomorphism,
then it may be decomposed as in the following commutative diagram
\begin{equation} \label{diag:isomorphism}
\begin{CD}
   G @>f>> H \\
   @VVV   @AAA \\
   G / \ker(f) @>>> f(G)
\end{CD}
\end{equation}
where $G / \ker(f)$ is the group $G$ with equality redefined to mean
$x =_{G / \ker(f)} y$ if and only if $f(x) =_H f(y)$,
where the downward arrow denotes the quotient map,
the lower horizontal arrow denotes the canonical group isomorphism,
and the upward arrow denotes the canonical injection.
All three of these maps are just $x \mapsto x$.
The quotient map and the canonical group isomorphism have equality defined
differently on their domain and codomain, whereas the canonical injection
has equality defined the same on both.

The theorem has zero interest in whether $f(G)$ and $H$ have some
sort of part-whole relationship (or exactly what kind if so).  We do have
$f(G)$ a subset of $H$ in the sense of ETCS because the canonical injection
is an injection.  But much more importantly, the fact that this arrow
indicates a group homomorphism says that $f(G)$ is a subgroup of $H$
rather than a mere subset.
Similarly $G / \ker(f)$ and $f(G)$ are isomorphic, hence
each is a subset of the other in the sense of ETCS.  And that is all group
theorists really need (or should want).

In diagram \eqref{diag:isomorphism} $\ker(f)$ is a subgroup of $G$ and
$f(G)$ is a subgroup of $H$, but $G / \ker(f)$ is not a subgroup of anything
in the diagram.  Rather, $G / \ker(f)$ is a factor group of $G$, that is,
$\ker(f) \times (G / \ker(f)) = G$.

So the notion of structure-preserving function (called homomorphism in group
theory) subsumes and generalizes the notion of subset or even
structure-preserving subset (here subgroup).
And this way of thinking is 100 years old and has gotten stronger and stronger
over the years.

\subsection{Constructive Analysis}

It may seem strange that we defined a quotient group in the preceding section
by redefining equality (rather than using equivalence classes), but this
is commonplace in constructive mathematics.  It goes back at least to
\citet{bishop} and is also seen in Martin-L\"{o}f type theory (MLTT)
\citep{martin-lof-mltt}, on which
homotopy type theory (HoTT)
\citep[hereinafter referred to as the HoTT book]{hott-book} is based.
\defcitealias{hott-book}{the HoTT book}

\Citet{bishop} defines a set to be rules for producing elements of the set
and rules for saying which elements of the set are equal.  So equality is
a defined relation, which must be an equivalence relation (reflexive,
symmetric, and transitive) and may be defined differently for each set.

In HoTT, if we want to be pedantically correct,
we write $=_A$ rather than just $=$ to mean equality for the type $A$.

So we did not need to think of equivalence classes at all (much less to
think of them as sets) to define a quotient group.

\subsection{Category Theory}
\label{sec:category}

We have already seen how in ETCS the set membership and subset relations are
defined in terms of functions.  In category theory, functions are generalized
to morphisms, which are like any abstraction: morphisms have some properties of
functions but not all properties.

\subsubsection{Definition and Axioms}

A category (like ETCS) has two sorts of variables, those ranging over
\emph{objects} and those ranging over \emph{morphisms}.
And it has a non-logical binary function
symbol $\circ$ for composition of morphisms.  And (according to taste, like
in ETCS) it has
a non-logical constant symbol $\idfun_A$ for every object $A$,
which are identity morphisms.

Every morphism $f : A \to B$ goes from one object to another.\footnote{So
if we are being very fussy, we can add to the formalism functions $\dom$
and $\cod$ each taking morphisms to objects such that $f : A \to B$ means
the same as $A = \dom(f)$ and $B = \cod(f)$.}
Like for functions, we can call $A$ the domain and $B$ the codomain of $f$.
But if we like to emphasize that morphisms are different from functions,
we can call $A$ the source and $B$ the target.
Morphisms are sometimes called arrows.

The axioms for category theory are very sparse.  For a diagram
\begin{equation} \label{diag:abc}
\begin{CD}
   A @>f>> B @>g>> C
\end{CD}
\end{equation}
there is a morphism $g \circ f : A \to C$.\footnote{That is,
$g \circ f$ exists if and only if $\dom(g) = \cod(f)$.}  And for a diagram
$$
\begin{CD}
   A @>f>> B @>g>> C @>h>> D
\end{CD}
$$
we have the associative law $h \circ (g \circ f) = (h \circ g) \circ f$.
And we also have the identity laws that say for any $f : A \to B$ we have
\begin{subequations}
\begin{align}
    f \circ \idfun_A & = f
    \label{eq:identity}
    \\
    \idfun_B \circ f & = f
    \label{eq:identity-too}
\end{align}
\end{subequations}
Identities are unique.\footnote{If $\idfun_A$ and $\idfun_A'$ were
both identites for object $A$, then the identity laws imply
$\idfun_A = \idfun_A \circ \idfun_A' = \idfun_A'$.\label{foot:identity}}

\subsubsection{Diagram Chasing}

Category theory makes much use of \emph{commutative diagrams}.
These are diagrams of directed graphs in which the nodes of the graph
are objects of the category and the edges of the graph are morphisms
of the category.  We say the diagram is commutative if parallel paths
are equal.

For example, to say that the following diagram is commutative
\begin{equation}
\label{diag:associative}
\begin{tikzcd}
    A \arrow{d}[swap]{f} \arrow{dr}{g \circ f} \\
    B \arrow{r}{g} \arrow{dr}[swap]{h \circ g} & C \arrow{d}[swap]{h} \\
    & D
\end{tikzcd}
\end{equation}
is to express the associative law: that all three paths from $A$ to $D$
are equal is to say that
$h \circ (g \circ f) = h \circ g \circ f = (h \circ g) \circ f$.

For another example, to say that the following diagram is commutative
\begin{equation}
\label{diag:associative}
\begin{tikzcd}
    A \arrow{r}{f} \arrow{d}[swap]{\idfun_A} \arrow{dr}{f} &
    B \arrow{d}{\idfun_B} \\
    A \arrow{r}{f} & B
\end{tikzcd}
\end{equation}
is to express the identity laws.

It is hard to describe how useful this concept is
without going through some complicated applications.
Commutative diagrams, like \eqref{diag:isomorphism} above,
organize whole proofs.  The diagram itself is (almost) the proof.
You just have to check that all the details it asserts are correct.

\subsubsection{Categories in Which Morphisms are Functions}

The axioms of ETCS say that ETCS is a category with sets as the objects
and functions as the morphisms.  But ETCS needs more axioms than just
the category axioms to make its
objects and morphisms act like sets and functions.

We note, if it isn't already obvious, that the objects of a category need
not be sets and do not have to have elements and that the morphisms of a
category need not be functions, and there need be no sense in which they
can be evaluated.  Examples below.

But for many categories, the objects are sets with extra structure and the
morphisms are structure-preserving functions.  Some of these are shown
in Table~\ref{tab:categories}.
\begin{table}
\caption{Categories of Structured Sets and Structure-Preserving Functions}
\begin{center}
\begin{tabular}{lp{2.75in}}
objects & morphisms \\
\hline
sets & functions \\
groups & group homomorphisms \\
rings & ring homomorphisms \\
fields & field homomorphisms \\
vector spaces & linear functions \\
topological spaces & continuous functions \\
metric spaces & uniformly continuous functions \\
topological groups & continuous group homomorphisms \\
smooth manifolds & differentiable functions \\
Lie groups & differentiable group homomorphisms \\
topological vector spaces & continuous linear functions \\
measure spaces & measurable functions \\
statistical models & functions measurable in the random variable
    and continuous in the parameter
\end{tabular}
\end{center}
\label{tab:categories}
\end{table}

\subsubsection{Isomorphisms}
\label{sec:iso}

In category theory, a morphism $f : A \to B$ is \emph{iso},
that is, $f$ is an \emph{isomorphism}, if there exists $g : B \to A$ such that
\begin{subequations}
\begin{align}
   g \circ f & = \idfun_A
   \label{eq:isodef-a}
   \\
   f \circ g & = \idfun_B
   \label{eq:isodef-b}
\end{align}
\end{subequations}
In all of the categories in Table~\ref{tab:categories} this category theory
notion of isomorphism coincides with the pre-existing notion of isomorphism
that existed before category theory was invented.

When $f$ is an isomorphism, we write a $g$ satisfying \eqref{eq:isodef-a}
and \eqref{eq:isodef-b} as $f^{-1}$ and say that $f^{-1}$ is the \emph{inverse}
of $f$.  This is OK because inverses are unique if they exist.\footnote{Suppose
\eqref{eq:isodef-a} and \eqref{eq:isodef-b} and also hold with $g$ replaced
by $g'$.
Then $g = g \circ \idfun_B = g \circ (f \circ g') = (g \circ f) \circ g'
= \idfun_A \circ g' = g'$.
\label{foot:inverse}}

\subsubsection{Preorders}
\label{sec:preorder}

A \emph{preorder} is a category in which for any two objects $A$ and $B$
there is at most one morphism $A \to B$.  Since there is one arrow $A \to A$,
the identity morphism, this is the only arrow $A \to A$.
Composition of morphisms says if there are arrows $A \to B \to C$, then there
is exactly one arrow $A \to C$.  Since the only issue is whether there is
an arrow between two objects, this looks like a relation (a binary predicate).
And the identity and composition laws say this is a reflexive and transitive
relation.  Such a relation is called a \emph{preorder}.
So to make this look more like conventional
mathematics we write $A \le B$ rather than $A \to B$.

Objects $A$ and $B$ in this category are isomorphic
if $A \le B$ and $B \le A$.\footnote{Proof: going back to arrow notation,
there can be only one arrow
$A \to B \to A$ or $B \to A \to B$, so these must be identity arrows;
hence $A \to B$ and $B \to A$ are inverse arrows.
This is essentially the same proof as in footnote~\ref{foot:terminal}.}
Orders do not allow this phenomenon.  So a preorder becomes a partial order if
we add the assumption that isomorphic objects are equal
($A \le B$ and $B \le A$ implies $A = B$).

And a partial order becomes a total order if we add the assumption that
for all objects $A$ and $B$, either $A \le B$ or $B \le A$.

The moral of the story is that morphisms in a category do not have to
be anything like functions, and categories do not have to be anything
like those in Table~\ref{tab:categories}

\subsubsection{Groupoids}
\label{sec:groupoid}

A \emph{groupoid} is a category in which every morphism is iso.
A \emph{group} is a groupoid with only one object.

If a group has only one object, call it $*$, how is that a group?
The morphisms of the category satisfy the group axioms.
Composition of morphisms is the group operation (called addition or
multiplication according to taste).  This operation is associative
by the associativity law for categories,
there is one identity morphism $\idfun_{\textstyle *}$
that behaves correctly because of the identity laws for categories.
Moreover, every morphism has an inverse by the assumption that all
morphisms are iso.

A group is thus a category in which the objects are uninteresting;
the arrows are everything.  But the arrows are nothing like functions.
This view of groups may seem strange, but was actually well known in
group theory before category theory was invented.
A \emph{permutation group} is exactly this construction when the object
is a finite set and the morphisms are the isomorphisms (invertible functions
from this finite set to itself).
An \emph{automorphism group} is exactly this construction when the object
is any mathematical object whatsoever and the morphisms are isomorphisms
in any sense that agrees with the definition in category theory:
in any category having an object $A$, the set of all isomorphisms $A \to A$
forms a group.

The idea of group presented in this section is different from the one
in Table~\ref{tab:categories}.  Here the category is just one group.
There it is all groups and the morphisms are group homomorphisms.
Here we seem to have lost the group homomorphisms, but we can get them
back (next section).

\subsubsection{Functors}

A \emph{functor} between categories is a function that maps
objects to objects and morphisms to morphisms and respects the laws,
that is, if $F : \mathcal{C} \to \mathcal{D}$ is the functor
and we have diagram \eqref{diag:abc}
in the source category $\mathcal{C}$ then we have the diagram
\begin{equation}
\label{diag:functor}
\begin{CD}
   F(A) @>{F(f)}>> F(B) @>{F(g)}>> F(C)
\end{CD}
\end{equation}
in the target category $\mathcal{D}$,
and the functor also maps identities to identities
$$
   F(\idfun_A) = \idfun_{F(A)}
$$

A functor between categories that are groups
(in the sense of Section~\ref{sec:groupoid})
is a group homomorphism in the sense of group theory.\footnote{Before
category theory, a group homomorphism was defined to be
a function $f : G \to H$ between groups that respected the group
multiplication $f(x y) = f(x) f(y)$ (if we denote the group operation by
multiplication).  But that is just what \eqref{diag:functor} says.
Then we have to prove that a group homomorphism respects identities.
Applying \eqref{diag:abc} with one or the other function being the
identity gives $\idfun_{\textstyle *} \circ f = f$
and $f \circ \idfun_{\textstyle *} = f$.  Applying \eqref{diag:functor} to
these gives $F(f) = F(f) \circ F(\idfun_{\textstyle *})$
and $F(f) = F(\idfun_{\textstyle *}) \circ F(f)$.  And this says
$F(\idfun_{\textstyle *})$ is an identity function.  Then uniqueness
of identities (Footnote~\ref{foot:identity}) finishes the proof.

We can also show that group homomorphisms preserve identities in the same
way.  Applying the functor laws to $f^{-1} \circ f = \idfun_A$
and $f \circ f^{-1} = \idfun_B$ gives
$F(f^{-1}) \circ F(f) = \idfun_{F(A)}$
and $F(f) \circ F(f^{-1}) = \idfun_{F(B)}$.  And this shows $F(f)$ is
an isomorphism.}

A functor between categories that are preorders (in the sense of
Section~\ref{sec:preorder} is a monotone function in the sense of the
theory of orders and preorders.\footnote{This is obvious.
Diagram \eqref{diag:functor}, translated into the symbology that replaces
$\to$ with $\le$ says $x \le y$ implies $f(x) \le f(y)$.}

\subsubsection{Homotopy}
\label{sec:homotopy}

Category theory was created \citep{eilenberg-maclane}
to do algebraic topology and homotopy theory.
What are they?

Algebraic topology is the study of topology using algebraic methods,
mostly using groups associated with topological spaces.
A homotopy is a continuous deformation of a continuous function.
Originally a homotopy $X \to Y$ was defined to be a (jointly) continuous
function $I \times X \to Y$ where $I$ denotes the unit interval
$\set{x \in \real : 0 \le x \le 1}$.
Here $I \times X$ denotes the product topological space, which has the product
topology.
The idea is that we have two functions $f_0 = f(0, \fatdot)$
and $f_1 = f(1, \fatdot)$
and more generally $f_t = f(t, \fatdot)$, and $f_0$ and $f_1$ are
continuous functions $X \to Y$ (morphisms of the category of topological
spaces and continuous functions) and $f_t$ changes continuously from $f_0$
to $f_1$ as $t$ goes from 0 to 1.

But this particular construction is irrelevant.  The actual concept is
continuous deformation.  When there is a homotopy like the one above we
say that $f_0$ and $f_1$ are \emph{homotopic} or \emph{homotopy equivalent}

This leads to the definition of \emph{homotopy category} in which the
objects are topological spaces and the morphisms are equivalence classes
of homotopy equivalent functions.
Alternatively, as in constructive mathematics, we can redefine what equals
means rather than use equivalence classes.  Morphisms are equal if they
are homotopic.

Isomorphisms in this category are the same as in any category.  A morphism
$f : A \to B$ is iso if there is a morphism $g : B \to A$ such that
\eqref{eq:isodef-a} and \eqref{eq:isodef-b} hold but now equals means
homotopic to, that is, $g \circ f$ is homotopy equivalent to the identity
function on $A$ and $f \circ g$ is homotopy equivalent to the identity
function on $B$.

One moral of the story is here is another category in which the morphisms
are not functions.

Now for the algebraic topology.  A \emph{path} in a topological space $A$
is a continuous function $I \to A$.  The relation of there being a path
from a point $x$ to a point $y$ is an equivalence relation (reflexive,
symmetric, and transitive).\footnote{Proof: the constant path takes $x$ to $x$,
hence reflexive.  Every path $p$ has an inverse $p^{-1}$ defined by
$p(t) = p(1 - t)$, $0 \le t \le 1$, hence symmetric.  Every two paths $p$
and $q$ such that the endpoint of $p$ is the startpoint of $q$ have a
path concatenation $r = p \cdot q$
defined by $r(t) = p(2 t)$, $0 \le t \le 1/2$
and $r(t) = q(2 t - 1)$, $1 / 2 \le t \le 1$, hence transitive.}

For any $x \in A$, we say the set of all paths from $x$ to $x$ is the
\emph{fundamental group} of $A$ for \emph{basepoint} $x$.  How is this
a group?  Path concatenation is the binary operation.  Every path from
$x$ to $x$ is then iso by the symmetry property.  This is the
notion of group-as-category as in Section~\ref{sec:groupoid}.

Two topological spaces that are topologically isomorphic (the term
for this used in topology before category theory was invented was
\emph{homeomorphic}) must have the same fundamental group but not
vice versa.  Thus the property of having the same fundamental group
is a weaker property than isomorphism.
Why would anyone want a weaker property?  Because it says something
different and perhaps useful and perhaps easier to calculate.

The trouble with equivalent concepts is the law of conservation of
mathematical difficulty.  Duality often does not help.
But a weaker concept perhaps can.

But this document is not about algebraic topology, so we leave the subject
here.

We do want to say something else about homotopy.
Since homotopies are continuous functions, we can make homotopies
of homotopies, jointly continuous functions $I \times I \times X \to Y$.
And homotopies of homotopies of homotopies,
jointly continuous functions $I \times I \times I \times X \to Y$.
And so on and on with arbitrarily many ``homotopies of.''
%%%%%%%%%% NEED FORWARD REFERENCE %%%%%%%%%%

\subsubsection{Category Theory as a Foundation of Mathematics}

\citet{lawvere-thesis,lawvere-cat-of-cat}
also proposed category theory as a foundation of mathematics.
The foundational theory is the category of all categories
in which the objects are categories and the morphisms are functors.

Of course, if the category of all categories had itself as one of the objects
we could reproduce Russell's paradox.
So like in NBG (von Neumann-Bernays-G\"{o}del) set theory, we do not allow
that.  We say that the objects of the category of all categories are
\emph{small categories} and that does not include the category of all
categories (which thus is analogous to the universe of NBG set theory,
which is a proper class).

Of course the category of all categories contains the category of sets and
functions.  Whether we take these to be ZFC sets and functions or NBG sets
and functions or ETCS sets and functions or HoTT sets and functions
\citepalias[Chapter~10]{hott-book} is a matter of taste.
But it also contains the rest of mathematics.

\subsubsection{Categorical Logic}

I preface this section with an admission that I am incompetent to write
on the subject, having never used it.  So I should not have the temerity
to write about it.  But I cannot just skip it either.  It is part of the
history.  But this section may have mistakes.

The simplest part of this story is that categorical logic generalizes
model theory (the branch of logic so called).  Theories are categories
(that's another part of the story) and models are functors from theories
into any other category.  There is no reason to have models in the category
of sets and functions (like classical model theory does).
In classical logic, models provide meaning (semantics).
Thus this feature of categorical logic is called \emph{functorial semantics}.

This is very powerful.  If one looks for models of the theory of groups in
the category of topological spaces and continuous functions, one discovers
topological groups.  If one looks in the category of smooth manifolds
and differentiable functions, one discovers Lie groups.

\subsection{Homotopy Type Theory}

For those brainwashed in logic and set theory, HoTT is really strange.
It has no need for logic and set theory.  It eats them.  It has logic
and set theory inside it.  They become parts of type theory.
So type theory is not a theory in first-order logic or any other logic.

\subsubsection{Types}

Types, the fundamental concept of type theory, are not defined
extensionally.
Instead types are intensional.  Type theory, according to
\citet{bauer}, is a theory of constructions.  Types are characterized by
rules for producing terms of the type.  This sounds like Bishop's definition
of set, but is more general.

In MLTT and HoTT the fundamental judgment is $a : A$
read as $a$ is a term of type $A$.
There are three different interpretations of $a : A$ used throughout HoTT
in addition to the ``neutral'' one just stated.
We can say $a$ is an element of the set $A$, or we can say $a$ is a point
in the space $A$ (which says the rules of $A$ give it additional structure
over and above being a mere set),
or we can say $A$ is a proposition (interpreted
as true if it is inhabited and false otherwise) and $a$ is a proof of $A$.

The latter shows us that type theory also eats proof theory.  To prove a
proposition $A$ is the same as producing a term of type $A$.  So we also
become completely uninterested in truth.  It is now defined to be types
being inhabited.

Some think it misleading when $p : P$ to call $p$ a proof of $P$.
Those people say $p$ is
a \emph{witness} to or \emph{evidence} of the truth of $P$.  But in MLTT
and HoTT the truth of $P$ does not mean anything other than that $p : P$ for
some $P$, that is, does not mean anything other than that $P$ is inhabited.

\subsubsection{Identity Types}

Also like in \citet{bishop}, as mentioned above, equality is a defined
relation for each type.
If $A$ is a type, and $x, y : A$, then we write $x =_A y$ to say that
$x$ and $y$ are equal terms \emph{according to the way equality is defined
for type $A$.}  But MLTT and HoTT go beyond that in considering this equality\
to be a type, denoted $(x =_A y)$ or $\text{Id}_A(x, y)$.  And, like any type,
this type can have terms.

So consider $p, q : (x =_A y)$.
What are $p$ and $q$?  Different (perhaps) ways $x$ and $y$ can be equal.
This caused a lot of confusion about MLTT until HoTT arose.
When we are thinking homotopically, we interpret $p$ and $q$ as paths
from $x$ to $y$ (as in algebraic topology and homotopy theory).  Then
it is no surprise at all that there can be more that one path from $x$
to $y$.  Similarly, if we are thinking that $x$ and $y$ are proofs of
$A$ considered as a proposition, it is no surprise that there can be
different proofs and maybe different ways proofs can be equivalent.
It is only when we say that there can be different ways that elements
of sets can be equal that this sounds bizarre to those brainwashed in
logic and ZFC-like set theory.

We should stress that, unlike in algebraic topology and homotopy theory,
these paths $p, q : (x =_A y)$ are not defined in terms of some other concepts.
Specifically they are not defined as functions $I \to A$ where $I$ is the
unit interval of the real line.  These are synthetic rather than analytic.
They are just always there for any types.  And they have the properties
given them for the rules of the type and no other properties.

Sometimes, like in the rest of math, we will omit the subscript on the
equals sign, writing $x = y$ instead of $x =_A y$, when the type
of $x$ and $y$ is obvious from the context.  But this does not indicate
a different concept.  In HoTT $x = y$ only makes sense when $x$ and $y$
have the same type.  If that type is $A$, then it means $x =_A y$.
In HoTT it makes no sense to ask whether terms of different types are equal.

Also, like in the rest of math, equality is an equivalence relation:
symmetric, reflexive, and transitive.  Reflexivity means $x = x$ for
any term $x$ of any type.  And what the preceding sentence means in HoTT
is that $(x = x)$ is an inhabited type (truth is inhabitation).  So for
any $x : A$ there exists
$$
   \refl_x : (x =_A x)
$$
which we can think of as the constant path from $x$ to $x$,\footnote{If
it were a function $f : I \to A$, where $I$ is the closed unit interval
of the real line, then it would be the function defined by $f(t) = x$,
$0 \le t \le 1$.} but, as always in HoTT, this is synthetic rather than
analytic, just there rather than defined in terms of something else.

Transitivity requires that $x, y, z : A$ and $x = y$ and $y = z$
imply $x = z$.  Propositions-as-types (Section~\ref{sec:curry-howard} below)
requires we write this as the function type (logical implication is a special
case of function)
$$
   ((x = y) \times (y = z)) \to (x = z)
$$
then the notion of currying says we can write a function of two variables
as a function of one variable whose value is a function of one variable,
that is, there is little difference between $f(x, y)$ and $f(x)(y)$, and
to be really pedantic (a set theorist is a person who thinks all functions
have only one argument) a function $f : A \times B \to C$ has arguments
that are ordered pairs so should be written $f\bigl((x, y)\bigr)$, which is
really ugly, so no one does that.  So the notation $f(x, y)$ is really
ambiguous about whether the type
is $A \times B \to C$ or $A \to B \to C$.\footnote{The latter means
$A \to (B \to C)$, the default being $\to$ is right associative.}
Thus we can rewrite the above as
$$
   (x = y) \to (y = z) \to (x = z)
$$
\citepalias[Lemma~2.1.2 and the recursor for product types,
equation (1.5.2)]{hott-book}.  This function of two variables, which are
paths $p : x = y$ and $q : y = z$ is called \emph{path concatentation}
and written $p \cdot q : x = z$.
Again, we can think of this the path that is obtained by following $p$
and then following $q$, but it is really synthetic rather than analytic:
transitivity requires it is just there.

Symmetry requires that $x, y : A$ and $x = y$ implies
$$
   (x = y) \to (y = x)
$$
and this function (implication) is written $p \mapsto p^{-1}$.
Again, we can think of $p^{-1}$ as $p$ travelled backwards,
but it is really synthetic rather than analytic.

It is a theorem that $\refl_x^{-1} \equiv \refl_x$
\citepalias[Lemma~2.1.1]{hott-book} ($\equiv$ means equal by definition
rather than just $=_A$ for some type $A$), which gives us another reason
for thinking of it as a constant path.  Also it acts like an identity
for path concatentation $p : x = y$ implies
$p = p \cdot \refl_y = \refl_x \cdot p$ and also all of the other identity
properties of a category
\citepalias[Lemma~2.1.4]{hott-book}.
Thus an identity type is a groupoid (Section~\ref{sec:groupoid} above)
when considered as a category
\citepalias[Section~2.1]{hott-book}.

All of the above may seem rather simple until we realize that any
statement about equality of paths involves the existence of paths between
paths (homotopies in homotopy theory, Section~\ref{sec:homotopy} above,
although, as usual, synthetic rather than analytic), that is, if
$p, q : (x =_A y)$, then we can wonder whether $p = q$ but that is
a question about whether the type $(p =_{(x =_A y)} q)$ is inhabited,
and that is a type of paths between paths (homotopies).
And then if that type has inhabitants $f$ and $g$, we can wonder whether
$f = g$ but then that is a question about whether the type
$f =_{(p =_{(x =_A y)} q)} g$ is inhabited, so a question about paths
between homotopies (homotopies of homotopies) and so forth ad infinitum
(just like in homotopy theory, Section~\ref{sec:homotopy} above).

\subsubsection{Propositions and Sets}

A type $A$ is a proposition \citepalias[Section~3.3]{hott-book} if it
has at most one path component, that is,
\begin{equation} \label{eq:proposition}
   \isProp(A) :\equiv \prod_{x, y : A} (x =_A y)
\end{equation}
where $:\equiv$ indicates a definition and the $\Pi$-type can be read
as either a Cartesian product or as logical for-all ($\forall$).
Here it is best read as the latter.  A type $A$ is a proposition
(\citetalias{hott-book} often says \emph{mere} proposition
to emphasize meaning this concept) if any two points are equal.
So in a sense it has only one point, although the homotopical interpretation
is only one path component (every two points are path connected).
This makes mere propositions quite boring.  The only thing you can say
about them is whether they are inhabited (true) or not (false).

A type $A$ is a set \citepalias[Section~3.1]{hott-book} if every path component
is contractible, that is,
\begin{equation} \label{eq:set}
   \isSet(A) :\equiv \prod_{x, y : A} \prod_{p, q : (x =_A y)}
   (p =_{(x =_A y)} q)
\end{equation}
This is somewhat confusing, especially the notation, $p =_{(x =_A y)} q$,
which is the identity type $=_B$ where $B$ is itself an identity type
$x =_A y$.

In short, a type is a set when it is homotopically boring.  Points
are either in the same path component or they aren't and the path
components are uninteresting (they are contractible).

This may not seem much like ZFC, but the analogy is actually very close.
These HoTT sets behave much like ETCS sets.  But ETCS has the axiom of choice
and HoTT doesn't (except if you add it as an extra assumption to prove
something), so the HoTT sets make a constructive set theory.

An element of a set is, as we said above, just $a : A$, a term of a type.
But a subset is weird: in one sense just like ZFC and in
another sense wildly different.  It is a $\Sigma$-type
\citepalias[Section~3.5]{hott-book}
\begin{equation} \label{eq:subset}
   \sum_{x : A} P(x)
\end{equation}
where each $P(x)$ is a mere proposition.  In another notation from the rest
of mathematics, this means the same thing as $\set{x \in A: P(x)}$.
So HoTT subsets are $\set{x \in A: P(x)}$.  We do not reify them as
anything ``more fundamental.''

This is like ETCS in that subsets are not ``in'' the sets they are subsets
of in any sense.  But otherwise it is very different.  In ETCS (recall)
subsets are injective functions.  In HoTT they are dependent pair types.

An element of the type \eqref{eq:subset} is a pair $(x, p)$
where $x : A$ and $p : P(x)$.  Of course, where there is no proof (that is,
when $P(x)$ is not inhabited), there is no pair ether.
So \eqref{eq:subset} consists of pairs $(x, p)$ such that $x$ is a point
of the subset and $p$ is a proof that this $x$ is in the subset.

You may think it inelegant that we are carrying the proofs along with
the points.  But think again.  The alternative is to search all of mathematics
for the proof when you need it to continue another proof.  By carrying it
along the computer proof assistant does not have to search for $p$ when it
is needed.  This is called \emph{proof-relevant} logic.

\subsubsection{The Curry-Howard Correspondence}
\label{sec:curry-howard}

MLTT and HoTT can seem horribly baroque compared to ZFC-like set theories.
So complicated!  They have lots of constructions and lots of rules
governing their constructions.  ZFC just has one non-logical symbol for the
membership relation and a few axioms.

But ZFC needs first order logic, and MLTT and HoTT do not.  So that is not
a fair comparison.  A good way to see that type theory is actually simpler
than logic plus set theory is to see that it does not need anything
not needed for them.  It just generalizes them.

So we copy Table~1 from \citetalias{hott-book} adding a bit from
Table~1.1 in \citet{rijke} to get our Table~\ref{tab:foo}
(p.~\pageref{tab:foo}).
\begin{table}
\begin{center}
\begin{tabular}{cllll}
Symbol & Types & Logic & Sets & Homotopy \\
\hline
$A$ & type & proposition & set & space \\
$a : A$ & term & proof & element & point \\
$B(x)$ & dependent type & predicate & family of sets & fibration \\
$b(x)$ & dependent term & conditional proof & family of elements & section \\
$0$ & empty type & $\bot$ & $\varnothing$ & empty space \\
$1$ & unit type & $\top$ & singleton set & one-point space \\
$A + B$ & sum type & $A \vee B$ & disjoint union & coproduct space \\
$A \times B$ & product type & $A \wedge B$ & set of pairs & product space \\
$A \to B$ & function type & $A \implies B$ & set of functions & function space
\\
$\sum_{x : A} B(x)$ & dependent pair type & $\exists_{x : A} B(x)$ &
    disjoint sum & total space \\
$\prod_{x : A} B(x)$ & dependent function type & $\forall_{x : A} B(x)$ &
    Cartesian product & space of sections \\
$\text{Id}_A$ or $=_A$ & identity type & equality (=) &
    $\set{(x, x) : x \in A}$ & path space \\
$A \to 0$ & & negation ($\neg A$) & &
\end{tabular}
\end{center}
\caption{Type Theory, Logic, Set Theory, and Homotopy Theory}
\label{tab:foo}
\end{table}

The last line of Table~\ref{tab:foo} does not really fit with the rest,
but we need it to fill out the Curry-Howard correspondence (how logic fits
into type theory).  If $A$ is a type, its negation
is defined to be the type $A \to 0$.   This function space is inhabited
if and only if $A$ is not inhabited.  In case $A \equiv 0$ the function
type $0 \to 0$ has exactly one term: the empty function having a vacuous
rule.\footnote{If $f : 0 \to 0$, then
can evaluate $f(x)$ for all $x : 0$ because there are no such $x$.
But there is such a function (all mathematicians agree).\label{foot:emptyfun}}
Conversely, in case $A$ is inhabited, there can be no function $A \to 0$
because there is an $x : A$ but there is no point in the empty type for $f(x)$
to be.

Another thing we should perhaps clarify in the table is that in set theory
$A \times B$ denotes the Cartesian product of $A$ and $B$, which is the
set of (ordered) pairs
$$
   \set{ (x, y) : x \in A \wedge y \in B }
$$
and, still in set theory, if $\mathcal{A} = \set{ A_i : i \in I }$ an
indexed family of sets then $\prod \mathcal{A} = \prod_{i \in I} A_i$
is the set of functions $f : I \to \bigcup \mathcal{A}$
satisfying $f(i) \in A_i$ for all $i$.  This is the general Cartesian
product of any family of sets.  One statement of the axiom of choice is
that if every element of $\mathcal{A}$ is nonempty, then $\prod \mathcal{A}$
is nonempty.  This includes the case where $\mathcal{A}$ is the empty family,
in which case the only inhabitant of $\prod \mathcal{A}$ is the identity
function on the empty set (the empty function of Footnote~\ref{foot:emptyfun}).

In HoTT ordered pair is a primitive concept.  We do not need an arbitrary
and bizarre coding into set theory such as the Kuratowski definition
$(a, b) = \{\{a\}, \{a, b\}\}$ where $\{a\}$ is the singleton set whose
only element is $a$ and $\{a, b\}$ is the unordered pair whose existence
is guaranteed by the pairing axiom of ZFC.\footnote{Of course, the existence
of singleton sets is also guaranteed by the pairing axiom plus the axiom of
extensionality because $\{a\} = \{a, a\}$.}  Then HoTT defines terms of
type $A \times B$ to be ordered pairs $(a, b)$ with $a : A$ and $b : B$.
This makes a non-dependent sum type $\sum_{a : A} B$ where $B$ does not
depend on $A$, the same as the product type $A \times B$.  So there is
some redundancy in the basic HoTT operations.
For the dependent function type, $f$ is a term of type $\prod_{x : A} B(x)$
if $f$ is a function such that $f(x) : B(x)$ for all $x : A$.  So here we
are just again reading the $\Pi$-type as logical for-all.

In HoTT \citepalias[Theorem~2.15.7]{hott-book} we do not need the axiom of
choice in order for a dependent product type to be inhabited.  It is just true.
However, this does not have the all of the logical consequences of the
classical axiom of choice.  (More on this after some preliminaries.)

\subsubsection{Propositional Truncation}

If $A$ is a type, then there is another type $\lVert A \rVert$,
called the \emph{propositional truncation} of $A$, defined
as if $a : A$ then $a : \lVert A \rVert$ (the rules for producing terms
of these types are the same) but if $a, b : \lVert A \rVert$, then $a = b$
(so equality is redefined so that all terms of the latter are equal).
This makes $\lVert A \rVert$ a mere proposition, regardless of what $A$ is.

In Table~\ref{tab:foo} we were a bit cavalier about two of the lines.
Various theorems in \citetalias{hott-book} prove that the operations take
mere propositions to mere propositions except for the sum types, which don't.
For example, the unit type $1$ is a mere proposition
but the sum type $1 + 1$ has two elements that
are not equal.  So it isn't a proposition.
So if we want to get what people usually
think of as logic, we have to squash the sums.  Logical or becomes
$\lVert A + B \rVert$ and the dependent sum type (logical exists) becomes
$\left\lVert \sum_{x : A} B(x) \right\rVert$
\citetalias[Definition~3.7.1]{hott-book}.

But \citetalias[Section~3.7]{hott-book}, also argues against the routine
use of propositional truncation.  If $P$, $Q$, and $R$ are mere propositions
and we want an implication $\lVert P + Q \rVert \to R$ it suffices to prove
$P + Q \to R$ because the recursion principle for propositional truncation
gives us the implication $\lVert P + Q \rVert \to R$.

What is the difference?  To prove $\lVert P + Q \rVert$, we prove $P$ or prove
$Q$ and forget which one.  To prove $P + Q$, we prove $P$ or prove
$Q$ and remember which one.  Mathematically, there is no
point to the forgetting.

And similarly for the dependent sum type, which we read as logical exists.
This is similar to the discussion of the subset relation \eqref{eq:subset}
above.
A term of $\sum_{x : A} P(x)$ is an ordered pair $(x, p)$ where
$x : A$ and $p$ is a proof of the proposition $P(x)$.  A term of its
propositional truncation forgets both $x$ and $p$ (more precisely,
for any $y : A$ and $q : P(y)$ we
have $(x, p) =_{\lVert \sum_{x : A} P(x) \rVert} (y, q)$.)
Mathematically, there is no point to the forgetting (propositional truncation).

So \citetalias{hott-book} never unnecessarily uses propositional truncation.
This gives its logic a different feel from the logic taught in philosophy
classes.

Now we can also clarify the difference between the so-called HoTT axiom
of choice mentioned above and the classical axiom of choice.  This is
\citepalias[Lemma~3.8.2]{hott-book}, if $A$ is a set and $B(x)$ is a set
for all $x : A$, then
$$
   \left( \prod_{x : A} \left\lVert B(x) \right\rVert \right) \to 
   \left\lVert \prod_{x : A} B(x) \right\rVert
$$
is read just like the classical axiom of choice: if each $B(x)$
is inhabited, then so is $\prod_{x : A} B(x)$.
But we repeat that without the propositional truncation this is just true
(indeed, here, trivially true, because when we remove the vertical bars
both sides of the arrow are the same, and any function type $A \to A$ is
inhabited by the identity function).

\subsubsection{Homotopy Equivalence}

The point of this section is that in HoTT isomorphism is too strong a concept.
Recall that the identity function $\idfun_A : A \to A$ maps $x \mapsto x$.
For $f : A \to B$ define
$$
   \qinv(f) :\equiv \sum_{g : B \to A}
   (f \circ g = \idfun_B) \times (g \circ f = \idfun_A)
$$
(the q in $\qinv$ is for quasi-inverse, the quasi being there to remind us
that $=$ may mean only equal up to homotopy, depending on how the types
are defined).
This looks like the HoTT analog of the regular definition of isomorphism
(we have to remember that $=$ means $=_{(A \to A)}$ or $=_{(B \to B)}$
as the case may be, a defined equivalence relation).

The trouble is, that (surprisingly) this is not a mere proposition
\citepalias[Theorem~4.1.3]{hott-book}.  So we could use propositional
truncation.  But (again perhaps surprisingly) the HoTT book does not.
Instead it defines three homotopically equivalent concepts, each of
which characterizes homotopy equivalence.  The simplest (IMHO) is
for $f : A \to B$
$$
   \binv(f) :\equiv \left( \sum_{g : B \to A} (f \circ g = \idfun_B)\right)
   \times \left( \sum_{h : B \to A} (h \circ f = \idfun_A)\right)
$$
This says that $f$ has a left inverse $h$ and a right inverse $g$ that are
not necessarily the same.
(The bi in $\binv$ is for bi-invertible, the property of having a left and
right inverse.)
We are used to inverses being unique if they exist.
They are in category theory.\footnote{Suppose $f$ has
a left inverse $h$ and a right inverse $g$, as in
the HoTT definition of $\binv(f)$.  Then
$$
   h = h \circ \idfun_B = h \circ (f \circ g)
   = (h \circ f) \circ g = \idfun_A g = g
$$
Hence $h = g$.  Moreover, this implies that any two-sided inverses must
be equal, as was proved in Footnote~\ref{foot:identity} above.}
and they are in set theory \citepalias{hott-book}.
For any function $f$, $\binv(f)$ is a mere proposition
\citepalias[4.3.2]{hott-book} and this is the definition that is the
key to the whole theory.

Now we define $\isEquiv(f)$ to be the same as $\binv(f)$.  And we
write $A \simeq B$ when $f : A \to B$ and $\isEquiv(f)$.

\subsubsection{Universes and Univalence}

In HoTT we allow types to have types.  This again raises the specter of
Russell's paradox.  So we collect types into universe types.  Then we can
consider a dependent type as a function $B : A \to \mathcal{U}$.
To be careful we cannot have $\mathcal{U} : \mathcal{U}$ because that
would allow Russell's paradox.  So we might have to have a sequence of
ever larger universes.  But normally, we do not need to worry about that
(and won't here).

Then we can consider for $A, B : \mathcal{U}$ the equality
type $(A =_\mathcal{U} B)$.  And it is provable in MLTT
\citepalias[Lemma~2.10.1]{hott-book} for $A, B : \mathcal{U}$ there
exists a function defined in the proof of the lemma
$$
   \text{idToEquiv} : (A =_\mathcal{U} B) \to (A \simeq B)
$$
Voevodsky proved that it is consistent to assume as an axiom that this
function is an equivalence, that is,
$$
   (A =_\mathcal{U} B) \simeq (A \simeq B)
$$
So this is the high water mark of the HoTT view of equivalence.  It
is sometimes misunderstood as saying that isomorphic types are equal,
but it is not that isomorphism is being reduced to equality (some sort
of skeletality condition) but rather than equality is being expanded to
be equivalent to equivalence.

We can now also clarify the status of the law of excluded middle (LEM) in
HoTT.  The following statement
$$
   \prod_{A : \mathcal{U}} \bigl(A + (A \to 0)\bigr)
$$
is incompatible with the univalence axiom
\citepalias[Theorem~3.2.2]{hott-book}.  So we cannot have that.

But we can consistently assume \citepalias[Section~3.4]{hott-book}
$$
   \prod_{A : \mathcal{U}} \Bigl(\isProp(A) \to \bigl(A + (A \to 0)\bigr)
   \Bigr)
$$
However, we do not routinely assume this (sometimes denoted $\text{LEM}_{-1}$)
because it makes our theorems not applicable to large parts of contemporary
mathematics.

As a simple example, notice that this document always says inhabited rather
than nonempty (= not empty, so inhabited = not not empty, and the classical
equivalence of these two concepts is literally double negation removal,
which is equivalent to LEM).  Often, just changing the way you say things
allows for constructive mathematics.

\subsubsection{Homotopy Levels}

This section is tricky and probably should be skipped at first reading.

Mathematicians see a pattern in the definitions of (mere) propositions
and (mere) sets in HoTT.

They start with contractibility.  A type $A$ is contractible
\citepalias[Section~3.11]{hott-book}
if it has exactly one path component, that is,
\begin{equation} \label{eq:contractible}
   \isContr(A) :\equiv \sum_{a : A} \prod_{b : A} (a =_A b)
\end{equation}
read there exists $a$ in $A$ such that for all $b$ in $A$ we have $a =_A b$.

When read homotopically, this says a bit more than path connected.
If $\isContr(A)$ is true, this is the same as it being inhabited,
which means it has an element, and an element of the dependent sum
type \eqref{eq:contractible}
is a pair $(a, f)$, where $a : A$ and $f : \prod_{b : A} (a =_A b)$,
and an element of this product type is a function $f : b \mapsto p$,
where $p : a =_A b$, which is a path from $a$ to $b$.
But the homotopic reading says this variation is continuous,
so $f$ is a homotopy that maps the constant path (one endpoint is always $a$)
to the identity path (the other endpoint ranges over $b : A$).
And in algebraic topology this is a stronger property than path connectedness.

Now for the pattern.  We define homotopy-$n$-types
\citepalias[Chapter~7]{hott-book}
by
\begin{align}
   \text{is-$(- 2)$-type}(A) & :\equiv \isContr(A)
   \\
   \text{is-$(n + 1)$-type}(A)
   & :\equiv
   \prod_{x, y : A} \text{is-$n$-type}(x =_A y),
   \qquad n \ge -2
   \label{eq:type-recursion}
\end{align}
The bizarre idea to start counting at $- 2$ is to match up with how
homotopy levels were defined in homotopy theory before HoTT was invented.

It is then a theorem
\citepalias[Lemma~???]{hott-book} that $\isProp(A)$ is equivalent (in a sense
to be discussed in Section~\ref{sec:equivalence} below)
to $\text{is-$(-1)$-type}(A)$.  Since both sides here are propositions,
this merely means there is a function (read as logical implication) from
each to the other.

\REVISED

so we generalize: we say that sets are homotopy zero types and
\citep[Chapter~7]{hott-book}
\MOVED[equation]
so by induction we define $n$-types for all natural numbers $n$,
except it is clear from the above that we can start lower.
If we let $(-2)$-types be contractible types
and $(-1)$-types be mere propositions, then the induction works
starting at $-2$.
This little bit of craziness is to match up with the way homotopy
levels are numbered in homotopy theory.

We need a proof that the first step of this recursion works,
that we go from contractibles to propositions using \eqref{eq:type-recursion}
in case $n = -2$, that is, we need to prove
\begin{equation} \label{eq:proposition-too}
\begin{split}
   \isProp(A)
   & \equiv
   \prod_{x, y : A} \isContr(x =_A y)
   \\
   & \equiv
   \prod_{x, y : A}
   \sum_{p : x =_A y} \prod_{b : x =_A y} (p =_{x =_A y} q)
\end{split}
\end{equation}
and we notice the right-hand side of \eqref{eq:proposition-too} is not
the same as the right-hand side of \eqref{eq:proposition}.
I don't even think \eqref{eq:proposition-too} is correct in the sense
that $\equiv$ means definitionally equal in HoTT and these may be only
homotopy equivalences \citet[Chapter~4]{hott-book}.  Or perhaps I should
say that we can define $\isProp$ the way \eqref{eq:proposition} does or
as \eqref{eq:proposition-too} does.

This is all very mysterious, so let us have some simplification.
\Citet[Lemma~3.11.3(iii)]{hott-book} says $\isContr(A)$ implies $A \simeq 1$,
where $\simeq$ is the symbol for homotopy equivalence.

\section{The Natural Numbers and Finiteness}

In HoTT the natural number type $\nats$

\begin{thebibliography}{}

\bibitem[Aluffi(2009)]{aluffi}
Aluffi, P. (2009).
\newblock \emph{Algebra, Chapter 0}.
\newblock American Mathematical Society, Providence, RI.

\bibitem[Bauer(2019)]{bauer}
Bauer, A. (2019).
\newblock Introduction to homotopy type theory.
\newblock A doctoral course on homotopy theory and homotopy type theory
    given by Andrej Bauer and Jaka Smrekar at the Faculty of Mathematics
    and Physics, University of Ljubljana, in the Spring of 2019.
\newblock The lectures on homotopy type theory by Bauer were recorded
    and are on \href{https://www.youtube.com/playlist?list=PL-47DDuiZOMCRDiXDZ1fI0TFLQgQqOdDa}{YouTube}.
\newblock The course notes for those lectures are on
    \href{https://github.com/andrejbauer/homotopy-type-theory-course}{GitHub}.

\bibitem[Birkhoff and Mac Lane(1941)]{birkhoff-maclane}
Birkhoff, G., and Mac Lane, S. (1941).
\newblock \emph{A Survey of Modern Algebra}.
\newblock Macmillan, New York.

\bibitem[Bishop(1967)]{bishop}
Bishop, E. (1967).
\newblock \emph{Foundations of Constructive Analysis}.
\newblock McGraw-Hill, New York.

\bibitem[Eilenberg and Mac Lane(1945)]{eilenberg-maclane}
Eilenberg, S., and Mac Lane, S. (1945).
\newblock General theory of natural equivalences.
\newblock \emph{Transactions of the American Mathematical Society},
    \textbf{58}, 231--294.
\newblock \doi{10.1090/S0002-9947-1945-0013131-6}.

\bibitem[Kleiner(2007)]{kleiner}
Kleiner, I. (2007).
\newblock \emph{A History of Abstract Algebra}.
\newblock Boston: Birkh\"{a}user.

\bibitem[Lawvere(1963)]{lawvere-thesis}
Lawvere, F. W. (1963).
\newblock Functorial Semantics of Algebraic Theories.
\newblock Ph.D. thesis, Columbia University.
\newblock Republished (2004) with an author’s comment and a supplement in
    \emph{Reprints in Theory and Applications of Categories}, \textbf{5},
    1--121.
\newblock \url{http://www.tac.mta.ca/tac/reprints/articles/5/tr5.pdf}.

\bibitem[Lawvere(1964)]{lawvere-etcs}
Lawvere, F. W. (1964).
\newblock An elementary theory of the category of sets.
\newblock \emph{Proceedings of the National Academy of Science of the U.S.A.},
    \textbf{52}, 1506--1511.
\newblock Expanded version (2005) with commentary by Colin McLarty and
    the author in \emph{Reprints in Theory and Applications of Categories},
    \textbf{11}, pp. 1--35.

\bibitem[Lawvere(1966)]{lawvere-cat-of-cat}
Lawvere, F. W. (1966).
\newblock The category of categories as a foundation for mathematics.
\newblock In S. Eilenberg, D. K. Harrison, S. MacLane, H. Röhrl (eds.),
    \emph{Proceedings of the Conference on Categorical Algebra - La Jolla
    1965}, 1--20.
\newblock Springer.
\newblock \doi{10.1007/978-3-642-99902-4}.

\bibitem[Leinster(2014)]{leinster}
Leinster, T. (2014).
\newblock Rethinking set theory.
\newblock \emph{American Mathematical Monthly}, \textbf{121}, 403--415.
\newblock \doi{10.4169/amer.math.monthly.121.05.403}.

\bibitem[Martin-L\"{o}f(1975)]{martin-lof-mltt}
Martin-L\"{o}f, P. (1975).
\newblock An intuitionistic theory of types: predicative part.
\newblock In: H. E. Rose, J. C. Shepherdson (eds.),
    \emph{Logic Colloquium ‘73, Proceedings of the Logic Colloquium,
    Studies in Logic and the Foundations of Mathematics}, \textbf{80}, 73--118.
\newblock Elsevier.
\newblock \doi{10.1016/S0049-237X(08)71945-1}.

\bibitem[Mazur(2008)]{mazur}
Mazur, B. (2008).
\newblock When is one thing equal to some other thing.
\newblock In \emph{Proof and Other Dilemmas: Mathematics and Philosophy},
B. Gold and R.~A. Simons (eds.)
\newblock Mathematical Association of America, Washington, DC.
\newblock pp.~221--242.

\bibitem[Rijke(2025)]{rijke}
Rijke, E. (2025).
\newblock \emph{Introduction to Homotopy Type Theory}.
\newblock Cambridge University Press.
\newblock \doi{10.1017/9781108933568}.
\newblock A preprint is available at
    \href{https://arxiv.org/abs/2212.11082}{Ar$\chi$iv}.

\bibitem[The Univalent Foundations Program(2013)]{hott-book}
The Univalent Foundations Program (2013).
\newblock Homotopy Type Theory: Univalent Foundations of Mathematics.
\newblock Institute for Advanced Study, Princeton, NJ.
\newblock \url{https://homotopytypetheory.org/book}.

\bibitem[van der Waerden(1930)]{van-der-waerden}
van der Waerden, B. L. (1930).
\newblock \emph{Moderne Algebra, Teil I} 1930, \emph{Teil II} 1931.
\newblock Springer-Verlag, Berlin.

\bibitem[Voevodsky(2015)]{voevodsky}
Voevodsky, V. (2015).
\newblock An experimental library of formalized mathematics based on
    univalent foundations.
\newblock \emph{Mathematical Structures in Computer Science}, \textbf{25},
    1278--1294.
\newblock \doi{10.1017/S0960129514000577}.

\end{thebibliography}

\end{document}

